\chapter*{Introducción}
\addcontentsline{toc}{chapter}{Introducción}
\markboth{Introducción}{Introducción}

El Método de los Elementos Finitos (MEF) constituye una de las herramientas centrales del análisis estructural y de la ingeniería computacional \autocite{zienkiewicz2013fem,cook2002concepts}, y su enseñanza forma parte del núcleo curricular de las carreras de ingeniería civil y mecánica \autocite[p.~1160]{perezsantiago2023fem}. Su comprensión exige articular tres planos que suelen presentarse de forma fragmentada al estudiante: la formulación matemática continua (elasticidad lineal, ecuaciones de equilibrio), su discretización mediante elementos isoparamétricos y cuadratura numérica, y la interpretación física de los resultados (campos de desplazamiento, deformación y tensión). EduFEM se propone como un software educativo de escritorio que integra estos tres planos en un único entorno interactivo, orientado a problemas bidimensionales de elasticidad lineal estática con elementos cuadriláteros Q4 y Q9.

\section*{Planteamiento del problema}

La enseñanza del MEF enfrenta una tensión recurrente entre el rigor formal y la intuición operativa. Por un lado, los textos clásicos desarrollan la teoría con un grado de abstracción considerable, donde las funciones de forma, la matriz Jacobiana, la matriz de deformación-desplazamiento y la integración de Gauss aparecen como pasos de una sucesión algebraica cuya conexión con la geometría concreta de un elemento no siempre resulta evidente para quien recién se inicia \autocite{bathe2014fem,reddy2006introduction}. Por otro lado, el software comercial de análisis por elementos finitos opera, desde la perspectiva del estudiante, como una caja negra: el usuario define geometría, material y cargas, y obtiene resultados, pero el procesamiento intermedio que transforma esas entradas en tensiones permanece oculto \autocite[p.~1160]{perezsantiago2023fem}.

Esta opacidad tiene consecuencias pedagógicas concretas. El estudiante puede aprender a operar una herramienta profesional sin comprender qué ocurre internamente, lo que dificulta diagnosticar resultados anómalos, evaluar la calidad de una malla o reconocer fenómenos numéricos como el bloqueo por cortante (\emph{shear-locking}) o los modos espurios (\emph{hourglass}). A esto se suma que las herramientas comerciales suelen ser costosas, no estar diseñadas con un propósito didáctico explícito \autocite[p.~1160]{perezsantiago2023fem} y ofrecer una localización al español limitada como recurso de aula, lo que restringe su utilidad para el aprendizaje. La brecha, en consecuencia, es doble: existe una dificultad de aprendizaje genuina vinculada a la naturaleza abstracta del método, y una insuficiencia de herramientas que la atiendan exponiendo el proceso de cálculo de manera transparente y guiada. Esa es la \emph{situación problemática} de la que parte el trabajo: en la investigación tecnológica el problema no se elige libremente, sino que resulta de interpretar y diagnosticar una realidad concreta \autocite[p.~83]{garciacordoba2005tecnologica}.

En términos formales, el \textbf{problema científico} de esta investigación se formula como la siguiente interrogante: ¿cómo debe construirse un software educativo para que exponga de forma transparente, interactiva y numéricamente verificable el canal de cálculo completo del MEF en elasticidad plana, de modo que cada etapa del procedimiento resulte observable y contrastable por el estudiante de ingeniería como apoyo a la comprensión de sus fundamentos? Es la única formulación del problema que el documento emplea: la \autoref{sec:matriz-consistencia} la retoma textualmente y la \autoref{sec:validacion-hipotesis} la responde. La variable causa es la forma en que la herramienta expone ese canal —transparencia, interactividad y verificabilidad numérica—; la variable efecto es la observabilidad y contrastabilidad del procedimiento intermedio que el software profesional encapsula. La insuficiencia de una herramienta con esos atributos es la contradicción que origina el trabajo \autocite[p.~7]{alvarez_metodologia}. El problema se delimita en lo disciplinar a la elasticidad lineal plana con elementos cuadriláteros Q4 y Q9 (véase \emph{Alcance y limitaciones}); en lo institucional, a la formación de grado en ingeniería civil, con la Carrera de Ingeniería Civil de la Universidad Autónoma Tomás Frías como destinataria inmediata de la herramienta; en lo espacial, al aula y al equipo personal del estudiante, sin requerir laboratorio de cómputo ni licencias; y en lo temporal, al desarrollo, verificación y validación del software realizados durante la gestión 2026.

El \textbf{objeto de estudio} \autocite[p.~8]{alvarez_metodologia} es el análisis estructural de medios continuos en elasticidad lineal plana por el Método de los Elementos Finitos: el canal de cálculo que va del mapeo isoparamétrico a la recuperación de tensiones con elementos cuadriláteros Q4 y Q9, sobre el que se construye el motor, se miden las variables y se contrastan los resultados. El \textbf{campo de acción}, la parte de ese objeto que el trabajo se propone mejorar \autocite[p.~13]{alvarez_metodologia}, es su enseñanza: el software educativo que expone ese canal de forma transparente, interactiva y verificable, y la memoria de cálculo que lo documenta sobre el modelo del propio estudiante.

\section*{Justificación}

El desarrollo de EduFEM se justifica en varios planos complementarios. En el plano académico, la herramienta busca cerrar la distancia entre la formulación teórica del MEF y su comportamiento numérico, haciendo visibles los pasos intermedios que el software profesional encapsula. La literatura sobre enseñanza del método reporta que la simulación interactiva y la manipulación directa de los objetos matemáticos del MEF conducen a una comprensión más intuitiva de sus procedimientos \autocites[p.~157]{lee2015interactive}[pp.~872, 874]{lee2015eigenmodes}, y que construir paso a paso las herramientas de cálculo junto con el alumno resulta más eficaz que la sola exposición teórica \autocite[p.~17]{bishay2020teaching}. Exponer el mapeo isoparamétrico, el Jacobiano, la matriz de deformación y la cuadratura de Gauss sobre el elemento real que el estudiante construyó permite anclar conceptos abstractos en una geometría tangible.

En el plano práctico, EduFEM está concebido como recurso de apoyo a la cátedra: es gratuito y de código abierto —su código fuente está disponible públicamente en línea\footnote{Repositorio del proyecto: \url{https://github.com/Sucullani/SoftwareED}}—, está íntegramente en español y organiza su interfaz según las tres fases canónicas del análisis (pre-proceso, proceso y post-proceso), lo que reproduce el flujo de trabajo profesional al tiempo que lo hace pedagógicamente legible. La generación automática de una memoria de cálculo —un documento que reconstruye el procedimiento completo, con sustitución numérica paso a paso, sobre los valores del modelo del propio usuario— ofrece un puente entre la exploración interactiva y el cálculo verificable a mano. Esta capacidad —que no se documenta entre los antecedentes de software educativo revisados (\autoref{sec:antecedentes})— convierte a la herramienta en un recurso de verificación: el estudiante puede reproducir cada operación con lápiz y papel o en una hoja de cálculo y contrastar sus resultados con los del programa, elemento por elemento.

En el plano metodológico, la decisión de construir un software propio en lugar de adoptar uno existente responde a la necesidad de controlar la transparencia del motor de cálculo y de alinear cada componente de la interfaz con un objetivo didáctico. Un motor desarrollado con bibliotecas numéricas de código abierto \autocite{harris2020numpy,virtanen2020scipy} permite exhibir las matrices y operaciones intermedias sin las restricciones de un producto cerrado, y verificar su precisión frente a referencias reconocidas.

\section*{Objetivos}

El objetivo general de este trabajo es desarrollar, en el lenguaje de programación Python y sobre bibliotecas numéricas de código abierto, EduFEM, un software educativo de escritorio para el análisis por el Método de los Elementos Finitos (MEF) de problemas bidimensionales de elasticidad lineal (tensión y deformación plana) con elementos cuadriláteros isoparamétricos Q4 y Q9, que integre la formulación matemática, la construcción y visualización interactiva del modelo y la verificación numérica, como apoyo a la enseñanza y el aprendizaje del MEF.

Los objetivos específicos que concretan este propósito son los siguientes:

\begin{enumerate}
    \item Fundamentar teóricamente el MEF para la elasticidad lineal plana con elementos isoparamétricos Q4 y Q9 —formulación variacional, mapeo isoparamétrico, integración numérica, ensamblaje, recuperación de tensiones, calidad de malla y criterios de verificación y validación—, como base común del motor de cálculo, de los módulos educativos y de la memoria de cálculo.
    \item Implementar un motor de cálculo MEF 2D con elementos isoparamétricos Q4 y Q9 (funciones de forma, Jacobiano, matriz B, matriz constitutiva D, rigidez por cuadratura de Gauss, ensamblaje, restricciones y recuperación de tensiones).
    \item Diseñar una interfaz gráfica interactiva organizada en pre-proceso, proceso y post-proceso, con un único lienzo compartido para construir, resolver y visualizar el modelo.
    \item Desarrollar módulos educativos que expongan, paso a paso y sobre el modelo real, las etapas del canal de cálculo del MEF desde el mapeo isoparamétrico hasta el ensamblaje global, complementados por el post-proceso y la memoria de cálculo para la solución del sistema y la recuperación de tensiones.
    \item Verificar y validar el software mediante el método de soluciones manufacturadas (MMS) y casos de referencia clásicos (viga de Timoshenko, membrana de Cook), contrastando contra soluciones analíticas y un software comercial (SAP2000).
    \item Generar la memoria de cálculo automática en PDF que documenta el canal de cálculo sobre el modelo del estudiante y, como complemento instrumental, ofrecer interoperabilidad de datos (DXF, CSV).
\end{enumerate}

\section*{Hipótesis y preguntas de investigación}

Por tratarse de un trabajo de desarrollo tecnológico antes que de un estudio experimental, este trabajo no plantea una hipótesis estadística, sino una \emph{hipótesis de diseño} y un conjunto de preguntas de investigación que orientan la construcción y la evaluación de la herramienta. La distinción no es una licencia del autor: en la investigación tecnológica la hipótesis es la solución tentativa a un problema concreto y su criterio de veracidad es la efectividad en la práctica, no la contrastación de un enunciado teórico \autocite[p.~85]{garciacordoba2005tecnologica}. La hipótesis de diseño se enuncia en forma comprobable: si se construye un software educativo que exponga el canal de cálculo completo del MEF en elasticidad plana de forma transparente —haciendo observable cada matriz y cada paso numérico—, interactiva y coherente con el flujo de trabajo profesional (pre-proceso, proceso y post-proceso), entonces (a) su motor reproducirá las tasas teóricas de convergencia del desplazamiento y acotará el error por debajo del $3\,\%$ en la flecha y del $1\,\%$ en la tensión normal frente a la solución analítica de la viga de Timoshenko y frente al modelo de SAP2000; (b) evidenciará sobre la membrana de Cook el bloqueo por cortante del elemento Q4 frente a la convergencia del Q9; y (c) cada etapa del canal resultará observable en un módulo educativo o en la memoria de cálculo. Cada cláusula puede resultar falsa —un motor con un error de orden, un Q4 que no bloqueara o una etapa sin exposición la refutarían— y los criterios con que se contrasta se fijan a priori en la \autoref{sec:validacion-resultados}. Que esa transparencia constituya, además, un apoyo al aprendizaje de los fundamentos del método es un \emph{supuesto} fundamentado en la literatura, que reporta una mejor comprensión con la simulación interactiva y reclama conectar los fundamentos con la operación del software \autocites[p.~157]{lee2015interactive}[p.~1168]{perezsantiago2023fem}; este trabajo no lo contrasta empíricamente, porque ese contraste excede su alcance, y lo deja planteado como recomendación. En la nomenclatura de tres variables del diseño clásico de investigación, el aporte es EduFEM; el atributo manipulado, la forma en que se expone el canal de cálculo —transparencia, interactividad y verificabilidad numérica—; y el efecto esperado, la observabilidad y contrastabilidad del procedimiento, en correspondencia con las variables del problema científico. La \autoref{sec:validacion-resultados} liga cada uno de esos tres elementos con los indicadores de la \autoref{tab:variables}.

De esta hipótesis se desprenden las siguientes preguntas de investigación que el trabajo procura responder:

\begin{itemize}
    \item \textbf{PI-1.} ¿Es posible construir un motor de cálculo MEF 2D con elementos Q4 y Q9 cuya precisión numérica sea verificable y comparable con soluciones analíticas y con software comercial de referencia?
    \item \textbf{PI-2.} ¿Qué organización de la interfaz y de los módulos educativos permite exponer el canal de cálculo del MEF (del mapeo isoparamétrico a la recuperación de tensiones) de manera transparente y coherente con el flujo de trabajo profesional?
    \item \textbf{PI-3.} ¿De qué modo la herramienta permite evidenciar fenómenos numéricos relevantes, como el bloqueo por cortante en elementos Q4 frente a la convergencia de los Q9?
\end{itemize}

\section*{Metodología de la investigación}

Por su finalidad, este trabajo es de carácter \emph{propositivo}: su núcleo es la propuesta y construcción de un artefacto (EduFEM) que resuelve un problema formativo concreto. Por su naturaleza se inscribe, además, en la \emph{investigación aplicada} de tipo \emph{tecnológico}: su producto es un artefacto de software cuya construcción y evaluación constituyen el núcleo del estudio. Por el abordaje del problema, el enfoque es \emph{cuantitativo}, pues la calidad del artefacto se establece midiendo tensiones, desplazamientos, errores numéricos y tasas de convergencia frente a referencias de exactitud conocida; los criterios de diseño pedagógico se fundamentan cualitativamente en la literatura sobre enseñanza del MEF \autocite[p.~168]{lee2015interactive}. Según su nivel u objetivo, la investigación es \emph{descriptiva} y \emph{comparativa}: caracteriza el comportamiento numérico del modelo y contrasta de forma sistemática el desempeño de los elementos Q4 y Q9, evidenciando fenómenos numéricos como el bloqueo por cortante del elemento Q4 frente a la convergencia del Q9.

El diseño metodológico —el modelo de simulación numérica que sostiene la investigación, las variables y su operacionalización, la matriz de consistencia, los casos de estudio y su criterio de selección, el procedimiento de verificación y validación con sus instrumentos, y los criterios de aceptación con que se contrasta la hipótesis— se desarrolla en detalle en la \autoref{sec:metodologia}.

\section*{Alcance y limitaciones}

El alcance de EduFEM está delimitado de manera explícita para garantizar profundidad sobre un dominio acotado. En este documento, la expresión \emph{análisis estructural} se emplea en su acepción restringida de análisis tenso-deformacional de elementos estructurales modelados como medios continuos bidimensionales mediante la teoría de la elasticidad lineal; no comprende el análisis de sistemas estructurales discretos —pórticos o reticulados—, ni la formulación de placas, cáscaras o problemas tridimensionales. La herramienta resuelve problemas \emph{bidimensionales} de \emph{elasticidad lineal estática}, en las dos idealizaciones clásicas de la mecánica de sólidos: \emph{tensión plana} (apropiada para cuerpos delgados cargados en su plano) y \emph{deformación plana} (apropiada para cuerpos prismáticos largos con sección y carga invariantes a lo largo del eje longitudinal) \autocite[p.~11]{timoshenko1970elasticity}. La discretización se realiza con \emph{elementos cuadriláteros isoparamétricos} de cuatro nodos (Q4) y de nueve nodos (Q9), que cubren respectivamente la interpolación bilineal y la cuadrática completa, y permiten ilustrar de manera contrastada los efectos del orden de interpolación sobre la precisión \autocites[pp.~137-138]{hughes2000fem}[pp.~164, 183-184]{onate2009structural}. El producto es un \emph{software de escritorio} de carácter \emph{educativo}, desarrollado íntegramente en \emph{español}, con una interfaz diseñada bajo criterios de minimalismo y claridad pedagógica.

La elección del dominio bidimensional no es solo una restricción de alcance, sino una decisión pedagógica fundamentada: el continuo plano ocupa el punto intermedio idóneo entre el análisis unidimensional de barras y el análisis tridimensional de sólidos. El elemento de barra, con el que suele iniciarse la enseñanza del método, resulta insuficiente para ilustrar los conceptos centrales de la formulación isoparamétrica —el mapeo entre coordenadas naturales y físicas, la matriz Jacobiana, la matriz de deformación-desplazamiento y la cuadratura en dos direcciones—, que solo emergen al pasar al continuo. El análisis tridimensional, en el otro extremo, sí contiene esos conceptos, pero su aparato algebraico los oscurece: la matriz constitutiva crece a $6\times6$, la de deformación-desplazamiento a $6\times24$ en un hexaedro de ocho nodos y la rigidez elemental a $24\times24$ \autocite[p.~218]{cook2002concepts}. Esas dimensiones vuelven impracticable el seguimiento paso a paso que constituye el núcleo didáctico de la herramienta. El plano es la mínima dimensión en la que aparece el verdadero carácter de continuo y, a la vez, conserva matrices de tamaño legible —$\bm{D}$ de $3\times3$ y $\bm{B}$ de $3\times8$ en el elemento Q4—, lo que permite exhibir cada operación de forma trazable. Además, la formulación isoparamétrica, la cuadratura de Gauss y el ensamblaje son directamente generalizables a tres dimensiones \autocite{bathe2014fem,zienkiewicz2013fem}. El dominio plano actúa así como un puente: lo que el estudiante comprende en dos dimensiones se extiende de manera natural al caso tridimensional, planteado como línea de continuación del proyecto.

Las limitaciones del trabajo se reconocen de manera complementaria. El software no aborda análisis tridimensional, comportamiento no lineal (material o geométrico), efectos dinámicos ni dependientes del tiempo, ni otros tipos de elementos (triangulares, de orden superior a Q9, o elementos especiales). El catálogo de materiales se restringe a comportamiento elástico, isótropo y lineal, definido por el módulo de elasticidad y el coeficiente de Poisson. La malla se construye manualmente —por tablas, por dibujo en el lienzo o por importación DXF— o por generación estructurada en los casos de referencia: el software no incorpora mallado automático. La validación numérica se apoya en casos con solución analítica conocida y casos de referencia reconocidos de la literatura, mientras que la evaluación empírica de la dimensión educativa con estudiantes queda fuera del alcance temporal del trabajo y se plantea como línea de continuación (véanse las recomendaciones). Estas fronteras responden a una decisión deliberada: privilegiar la transparencia y la solidez de un núcleo conceptual bien acotado por sobre la cobertura amplia de funcionalidades, coherente con el propósito didáctico de la herramienta.

\section*{Estructura del documento}

El presente documento se organiza en tres capítulos, además de esta introducción, de un capítulo final de conclusiones y de un conjunto de anexos. El \autoref{cap:marco_teorico} desarrolla el marco teórico: abre con los antecedentes y el estado del arte del software para la enseñanza del MEF —que sitúan el trabajo y justifican su oportunidad— y continúa con los fundamentos del método aplicados a la elasticidad lineal bidimensional (formulación variacional, funciones de forma, mapeo isoparamétrico, matriz Jacobiana, matriz de deformación, matriz constitutiva, cuadratura de Gauss, ensamblaje y recuperación de tensiones), la calidad de malla y los criterios de verificación y validación. El \autoref{cap:diseno} expone el diseño e implementación del modelo: abre con el diseño metodológico —el tipo de investigación y su modelo de simulación numérica, las variables y su matriz de consistencia, los casos de estudio, el procedimiento de verificación y validación, y los criterios de aceptación con que se contrasta la hipótesis— y continúa con el diseño e implementación de la herramienta —los requisitos y decisiones de diseño, el stack tecnológico, la arquitectura en torno al modelo de proyecto, el motor de cálculo, el pre-proceso interactivo, los módulos educativos, el post-proceso y la generación de la memoria de cálculo—. El \autoref{cap:resultados} presenta y analiza los resultados: expone el producto final, mide los datos, ejecuta la verificación y validación numérica contrastando EduFEM frente a soluciones analíticas y software comercial, interpreta los hallazgos y contrasta la hipótesis. Los objetivos específicos se reparten entre los capítulos: el primero se cumple en el \autoref{cap:marco_teorico}; el segundo, el tercero, el cuarto y el sexto se materializan en el \autoref{cap:diseno}; y el quinto se cumple en el \autoref{cap:resultados}, que reporta además la verificación del motor, la cobertura de los módulos y la prueba de interoperabilidad. A continuación, el capítulo de conclusiones responde a cada objetivo en ese mismo orden, señalando el capítulo que lo desarrolla, discute el cumplimiento de la hipótesis de diseño y agrupa las recomendaciones por capítulo. Sigue la lista de referencias bibliográficas. Las citas y esa lista siguen la norma Vancouver (ICMJE y NLM), con numeración arábiga por orden de primera mención —entre corchetes en lugar de superíndice, para poder acompañar el número con la página citada—, títulos de revista completos y conjunción antes del último autor; el localizador de página acompaña a toda cita que sostiene una ecuación, un valor numérico, un umbral, una definición o una atribución concreta, y se omite en las citas de encuadre. La presentación del documento sigue la séptima edición de las normas APA. Por último, los anexos reúnen el manual de instalación (\autoref{anexo:instalacion}), el manual de uso (\autoref{anexo:manual}), una selección de listados de código del motor numérico (\autoref{anexo:codigo}), los resultados extendidos de verificación y validación (\autoref{anexo:vyv}), el modelo de validación en SAP2000 y la solución analítica de Timoshenko (\autoref{anexo:sap2000}), los formatos de interoperabilidad (\autoref{anexo:formatos}) y una memoria de cálculo resuelta paso a paso sobre el ejemplo canónico (\autoref{anexo:memoria}).